\documentclass[11pt]{article}
\usepackage{amsmath,amssymb,amsthm,hyperref}
\newtheorem{theorem}{Theorem}
\newtheorem{lemma}[theorem]{Lemma}
\newcommand{\E}{\mathbb E}
\newcommand{\Prob}{\mathbb P}
\title{Uniform approximation by independently relabelled palindrome blocks}
\author{Research note; independently reviewed by two other agents}
\date{30 September 2026}
\begin{document}
\maketitle

Let $\rho_1,\ldots,\rho_m$ be independent uniform permutations of $[n]$.
Concatenate the words $\rho_i\rho_i^{\mathrm{rev}}$.
Each die has $2m$ equally likely faces. For an output order $\sigma\in S_n$,
let $R_\sigma=n!P(\sigma)$ be its likelihood ratio relative to uniform.

\begin{theorem}\label{thm:main}
For $n\geq3$, $0<\varepsilon\leq1$, put
\[
 p=\max\left(2,\frac{\log(2n!)}{\log3}\right),\qquad
 m\geq\max\left(n,
       \left(\frac{4n^3\sqrt{p}}{\varepsilon}\right)^{2/5}\right).
\]
With probability at least $1/2$ over the independent block permutations,
\[
 |R_\sigma-1|\leq\varepsilon\qquad(\sigma\in S_n).
\]
Consequently $O(n^{7/5}(\log n)^{1/5}\varepsilon^{-2/5})$ faces per die
suffice for pointwise relative $\varepsilon$-fairness.
\end{theorem}

\begin{lemma}[A bounded canonical-kernel estimate]\label{lem:chaos}
Let $X_1,\ldots,X_m$ be independent. Suppose
\[
 F=\sum_{\varnothing\ne I\subseteq[m]}F_I(X_I)
\]
is a finite sum, with each $F_I$ centered separately in each of its arguments,
and $\|F_I\|_\infty\leq\beta^{|I|}$. For real $p\geq2$, put
$a=2\beta\sqrt{m(p-1)}$. If $a<1$, then
\[
 \|F\|_p\leq\sum_{\ell\geq1}\frac{a^\ell}{\sqrt{\ell!}}
             \leq\frac a{1-a}.
\]
\end{lemma}
\begin{proof}
Write $F^{(\ell)}=\sum_{|I|=\ell}F_I$ and introduce independent copies $X_i'$.
Let $\Delta_I F_I$ be the iterated difference in all coordinates of $I$,
using $X_i-X_i'$ in the usual difference-operator notation.
Separate centering gives
$\E_{X'}\sum_{|I|=\ell}\Delta_I F_I=F^{(\ell)}(X)$.
Jensen's inequality bounds its $L^p$ norm by that of the differenced sum.
Swapping each pair $(X_i,X_i')$ independently leaves its distribution
unchanged and multiplies $\Delta_I F_I$ by the product of the corresponding
Rademacher signs. Conditional on the pairs, the Rademacher
Bonami--Beckner inequality \cite{bonami,odonnell} for a homogeneous degree-$\ell$ polynomial gives
\[
 \left\|\sum_{|I|=\ell}\Delta_I F_I\right\|_p
 \leq (p-1)^{\ell/2}
       \left(\sum_{|I|=\ell}(2^\ell\beta^\ell)^2\right)^{1/2}
 \leq \frac{(2\beta\sqrt{m(p-1)})^\ell}{\sqrt{\ell!}}.
\]
The estimate is uniform in the conditioning. Sum over $\ell$ using
Minkowski's inequality.
\end{proof}

\begin{proof}[Proof of Theorem~\ref{thm:main}]
Fix $\sigma$. Let $L=(L_1,\ldots,L_m)$ have the multinomial distribution
with $n$ trials and equal cell probabilities $1/m$. Reading $\sigma$ from
left to right, split it into consecutive segments $A_i$ of lengths $L_i$.
Let $q_i(A_i)$ be the conditional probability of the prescribed internal
order within the $i$th palindrome block, and put
\[
 r_i(A_i)=L_i!q_i(A_i),\qquad \delta_i(A_i)=r_i(A_i)-1.
\]
For empty segments set $r_i=1$. Directly conditioning on block choices gives
\[
 R_\sigma=\E_L\prod_{i=1}^m r_i(A_i).
\]
Uniform random relabelling gives $\E_{\rho_i}\delta_i(A)=0$ for every fixed
segment $A$. Furthermore $\delta_i(A)=0$ for $|A|\leq2$, since each
palindrome block is pairwise fair.

A prescribed order of $k$ distinct labels has at most two subsequence
occurrences in one palindrome, so $q_i(A)\leq2^{1-k}$.
For $k\geq3$ this implies
\[
 |\delta_i(A)|\leq k!2^{1-k}.
\]
Expanding the product gives
\[
 R_\sigma-1=\sum_{\varnothing\ne I\subseteq[m]}F_I,
 \qquad F_I=\E_L\prod_{i\in I}\delta_i(A_i).
\]
These kernels are centered separately in every block permutation indexed
by $I$. If $|I|=\ell$, their uniform bound is
\begin{align*}
 |F_I|
 &\leq\sum_{\substack{k_i\geq3\ (i\in I)\\d=\sum_{i\in I}k_i\leq n}}
 \frac{n^{\underline d}}{m^d}
       (1-\ell/m)^{n-d}\prod_{i\in I}2^{1-k_i}\\
 &\leq\left(\sum_{k\geq3}2^{1-k}(n/m)^k\right)^\ell
 =\beta^\ell,
 \qquad \beta=\frac{n^3}{4m^3(1-n/(2m))}.
\end{align*}
The first line is just the selected-cell multinomial marginal; its
$k_i!$ denominators cancel the factorials in the bound for $\delta_i$.
All omitted terms in the next line are nonnegative.

Since $m\geq n$, we have $\beta\leq n^3/(2m^3)$ and therefore
\[
 a=2\beta\sqrt{m(p-1)}\leq\frac{n^3\sqrt p}{m^{5/2}}
 \leq\frac\varepsilon4.
\]
Lemma~\ref{lem:chaos} implies $\|R_\sigma-1\|_p\leq\varepsilon/3$.
Markov's inequality gives
$\Prob(|R_\sigma-1|>\varepsilon)\leq3^{-p}$.
A union bound over $n!$ orders proves the theorem.
\end{proof}

\paragraph{What this does and does not establish.}
This argument controls the complete nonlinear Hoeffding expansion. It does
not use a first-order approximation. It does not establish the sharper
$n^{13/10}$ heuristic suggested by internal dependence among overlapping
triples. No claim of optimality or literature novelty is made here.
\begin{thebibliography}{9}
\bibitem{bonami} Aline Bonami, \emph{\'Etude des coefficients de Fourier des fonctions de $L^p(G)$}, Annales de l'Institut Fourier \textbf{20} (1970), no.~2, 335--402, \href{https://aif.centre-mersenne.org/item/AIF_1970__20_2_335_0/}{doi:10.5802/aif.357}.
\bibitem{odonnell} Ryan O'Donnell, \emph{Analysis of Boolean Functions}, Cambridge University Press, 2014, Theorem~9.21; \href{https://www.cs.cmu.edu/~odonnell/papers/Analysis-of-Boolean-Functions-by-Ryan-ODonnell.pdf}{author's complete book manuscript}.
\end{thebibliography}
\end{document}
